The math for this project will be split into three parts:
\begin{enumerate}
	\item Calculating the survivor fraction
	\item coordinate randomization
	\item mass distribution 
\end{enumerate}
\subsection{calculating the survivor fraction}
The survivor fraction is how we measure whether the Uranium has reached the critical mass. It is calculated by dividing the previous "generation" of fission by the current "generation" of fission. 

If the survivor fraction ($f$) is lower than 1, we can conclude that we are below the critical mass, meaning that the reaction will die out.

If the survivor fraction is higher than 1, we can conclude that the reaction has gone supercritical.

By changing the mass (user input), we can try to find the point at which the survivor fraction is 1. To make sure this is accurate, we will allow the user to insert an amount of points ($N$) that need to be simulated. We will then simulate those $N$ random points inside the volume. this will be our first generation. We will then simulate fission for each of those $N$ points. If we check whether the neutrons are inside the volume, we can calculate the survivor fraction with the formula:
\begin{equation}
	f=\frac{N_{\text{inside volume}}}{N_{\text{gen 1}}}
\end{equation}
To be able to check whether the neutrons are inside the volume, we first need to build a randomizer for the movement of the neutrons inside the simulations.
\subsection{coordinate randomization}\label{randomizer}
Using a randomizer to generate random uniformly distributed points, we can calculate the path of the neutrons. If the neutrons have starting coordinates ($x_0$, $y_0$, $z_0$), we can describe their endpoints ($x_1$, $y_1$, $z_1$) using the equations:
\begin{equation}
	\begin{cases}
		x_1= x_0+d\;\sin\theta\;\cos\phi \\
		y_1=y_0+d\;\sin\theta\;\sin\phi \\
		z_1=z_0+d\;\cos\theta
	\end{cases} \label{newcoords}
\end{equation}
Here, $d$ is the distance the neutron travels and $\theta$ and $\phi$ are the angles of spherical coordinates. We will need to randomize these three variables.

$d$ Is a random variable between 1 and 0. $\phi$ Is a random variable between $2\pi$ and 0. Both of these variables are uniformly distributed.

$\theta$ is a bit trickier, for $\cos\theta$ is the uniformly distributed variable. We will thus have to randomize a value between $-1$ and $1$ and use the $\arccos$ function to  calculate the angle $\theta$.
\subsection{mass distribution}
For the mass distribution, we have to take into account that we have two possible shapes to simulate: a beam (or cube) and a sphere.
\subsubsection{beam} \label{beam}
For the beam, we will assume that the mass ($M$) is equal to the volume ($V$). We can then change the shape of the beam using a scale factor ($S$). The influence of the scale factor can be shown with the following equations\supcite{opgave}:
\[M=V=a^2b=\frac{a^2}{S}=b^3S^2\]
Here, $a$ is the width and the height of the beam and $b$ is the depth of the beam. If we want to calculate $a$ and $b$, we can derive following equations\supcite{opgave}:
\begin{align}
	a&=\sqrt[3]{M*S} \\
	b&=\sqrt[3]{\frac{M}{S^2}} \label{depth}
\end{align}
We can thus see that for $S=1$, we get a perfect square.

We will assume that this beam is centered around (0, 0, 0). We can then create a starting point randomizer, using a uniform distribution between the following endpoints:
\begin{displaymath}
	\begin{cases}
		x&\in\left[-\frac{a}{2}, \frac{a}{2}\right] \\
		y&\in\left[-\frac{a}{2}, \frac{a}{2}\right] \\
		z&\in\left[-\frac{b}{2}, \frac{b}{2}\right] \\
	\end{cases}
\end{displaymath}
We can also use these conditions for the evaluation to check whether a point is inside the beam object.
\subsubsection{sphere} \label{sphere}
For the sphere, we will also assume that the mass is equal to the volume. For this class, we don't need to implement a scale factor. We do however need to calculate the radius. this can be done by using the formula:
\[V=\frac{4}{3}\pi\;r^3\]
If we then substitute $V$ for $M$, we can derive the formula for the radius:
\begin{equation}
	r=\sqrt[3]{\frac{3}{4}\;M}
\end{equation}

For the random starting points inside the sphere, we can reuse the code for coordinates of a neutron. Because the distance for the neutrons is between 0 and 1, we can simply multiply this number with the radius to get the desired output.

To evaluate whether a point is inside this sphere, we calculate the distance from the center of the circle (0, 0, 0) and compare this to the radius.